\scanpage{255}

% folio and opening continuation
\mspage{8}à isom unique près. En fait, quand

% derivator defined by a pair
\[\mathbb D=\mathbb D_{\deleted{\unsure{D|M}}\mathcal M,\Sigma}\]

% between first two displays
est défini par une catégorie de modèles $(\mathcal M,\Sigma)$, alors il y a une définition naturelle \deleted{de l'espèce} d'une opération $\mathrm{Hom\,ext}$ (Hom ``externe'') à valeurs en $\mathrm{Hot}$, \deleted{\unsure{dans $\mathcal A$}} savoir un foncteur

% display (11), two lines with a left brace
\[
\left\{\begin{aligned}
\mathcal A^{\unsure{\circ|0}}\times\mathcal A&\longrightarrow\mathrm{Hot}\\
(M,N)&\longmapsto\mathrm{Hom\,ext}(M,N)
\end{aligned}\right.\tag*{(11)}
\]

% after (11), fundamental category
où $\mathcal A=\mathcal M\Sigma^{-1}$ est la cat. fond. de $\mathbb D=\mathbb D_{\mathcal M,\Sigma}$.

% middle paragraph, verification and source
\deleted{\unsure{Quand} \illegible{1}} Il est vrai que je n'ai pas vérifié que \deleted{\unsure{si|a}}si p.ex. $(\mathcal M,\Sigma)$ est sous-jacente à une catégorie de modèles \unsure{quillénienne|quilléniennes}, de sorte que (par \unsure{Thomason}) $\mathbb D$ est bien un dérivateur, que cette notion de $\mathrm{hom\,ext}$ \deleted{\unsure{obéit à}}\inserted{\unsure{satisfait|satisfais}} bien à une formule de la forme (10).

% lower paragraph, additional datum
\uline{Ici} je vais la poser, \deleted{\unsure{ou}} \unsure{disons|décrire}, en donnée \deleted{\unsure{et axiome}} supplémentaire : donnée de (11), et d'un iso. fonctoriel

% (12), three expressions at left joined by vertical double strokes
\[
\begin{aligned}
\mathrm{Hom}_{\unsure{\xi|\zeta}}(M,N)&\simeq\mathrm{Hom}_{\mathrm{Hot}}(\xi,\mathrm{Hom\,ext}(M,N))\\
\mathrel{\Vert}\hspace{1em}&\\
\mathrm{Hom}(\xi\times_{\mathbb D}M,N)&\\
\mathrel{\Vert}\hspace{1em}&\\
\mathrm{Hom}(M,\mathrm{hom}_{\mathbb D}(\xi,N))&
\end{aligned}\tag*{(12)}
\]

% archive sheet number: 255
